

\section{拍与拍频}

\begin{definition}[拍与拍频]
当两个简谐运动的方向相同，但频率存在微小差异时，合成后的运动将不再是单纯的简谐运动，而是会呈现出一种振幅随时间作周期性缓变的有节奏的强弱交替现象——这就是波动学中著名的拍（Beats）现象。
\[x_1=A\cos(2\pi\nu_1 t),x_2=A\cos(2\pi\nu_2 t)\implies x=2A\cos\left[\pi(\nu_2-\nu_1)t\right]\cos\left[2\pi\frac{\nu_1+\nu_2}{2}t\right]\]
方程尾部 $\cos 2\pi \frac{\nu_2 + \nu_1}{2} t$ 代表了高频的快速振动。它决定了质点往复运动的实际快慢。故合振动主频率$$\nu = \frac{\nu_1 + \nu_2}{2}$$
由于 $|\nu_2 - \nu_1| \ll \nu_1 + \nu_2$，方程前部的括号项 $\cos 2\pi \frac{\nu_2 - \nu_1}{2} t$ 随时间变化得极其缓慢。我们可以将这部分视作整体，定义为合运动的瞬时视在振幅 $A(t)$：$$A(t) = \left| 2A_1 \cos 2\pi \frac{\nu_2 - \nu_1}{2} t \right|$$，这个变化的振幅完成了一次“从强到弱、再从弱到强”就是一个周期（拍的周期） ，所以$T_{\text{拍}}$ 严格等于该项余弦函数周期的二分之一：$$T_{\text{拍}} = \frac{1}{2} \cdot \left( \frac{1}{\frac{\nu_2 - \nu_1}{2}} \right) = \boxed{\frac{1}{\nu_2 - \nu_1}}$$根据频率与周期互为倒数的定义，合振动强弱交替的瞬时频率（即拍频）为：$$\nu_{\text{拍}} = \frac{1}{T_{\text{拍}}} = \boxed{|\nu_2 - \nu_1|}$$
\end{definition}

\begin{note}
在平面直角坐标系内，引入两个分简谐运动对应的旋转矢量 $\vec{A}_1$ 和 $\vec{A}_2$：由于两者的角速度不相等（$\omega_2 > \omega_1$）。如果我们站在矢量 $\vec{A}_1$ 的视角上看（以 $\vec{A}_1$ 为参考系），矢量 $\vec{A}_2$ 就会以极慢的相对角速度 $\Delta \omega = \omega_2 - \omega_1$ 顺时针或逆时针悄悄转圈。
\begin{itemize}
\item 同向重合时（相位差 $\Delta \phi = 2k\pi$）：两矢量首尾相连、并向延伸，此时合振动振幅 $\vec{A}$ 叠加达到极大值 $2A_1$（相长干涉，声音最响）。
\item 反向重合时（相位差 $\Delta \phi = (2k+1)\pi$）：两矢量背道而驰、彼此拉跨，此时合振动振幅 $\vec{A}$ 抵消达到极小值 $0$（相消干涉，声音最弱）。
\end{itemize}
由于 $\vec{A}_2$ 相对于 $\vec{A}_1$ 旋转一周（转过 $2\pi$ 弧度）就对应完成一次“强-弱-强”的拍流转，其对应的拍的周期自然为：$$T_{\text{拍}} = \frac{2\pi}{\omega_2 - \omega_1} = \frac{2\pi}{2\pi \nu_2 - 2\pi \nu_1} = \frac{1}{\nu_2 - \nu_1}$$对应的拍频自然收拢为：$$\nu_{\text{拍}} = \nu_2 - \nu_1$$
\end{note}

\begin{exercise}[拍频听起来在说什么]
若两音叉频率分别为 \(256\ \text{Hz}\) 和 \(258\ \text{Hz}\)，那么你每秒会听到多少次明显的强弱变化？
\end{exercise}

\begin{solution}
拍频讨论的是“强弱起伏有多快”，它等于两个频率之差：
\(\nu_{\text{拍}}=\abs{258-256}=2\ \text{Hz}\)，而主频约为 \((256+258)/2=257\ \text{Hz}\)。这表示每秒会听到 2 次明显的强弱变化；但真正声音本身的振动并不慢，它的主频仍接近 \(257\ \text{Hz}\)。
所以“每秒听到 2 次忽强忽弱”并不等于“声音本身的频率只有 2 Hz”。拍频题里一定要先把“振动频率”和“强弱变化频率”分开看。
\end{solution}
